\section{Requirements and Data Handling}
\subsection{Functional Baseline}
The requirements below summarize the identifiers maintained in
\texttt{docs/requirements.md}. Acceptance is based on observable behavior.

\begin{center}
\renewcommand{\arraystretch}{1.2}
\begin{tabularx}{\textwidth}{l X}
\hline
\textbf{IDs} & \textbf{Required behavior} \\
\hline
R01--R02 & Validate a local sidebar upload, extract visible names, and allow
correction, removal, and addition before explicit approval. \\
R03--R04 & Preserve a stable identity for every entry, including identical
display names, and generate one output row per approved entry. \\
R05--R07 & Create a distinct template copy for every entry, preserve questions
and independent responses, and write the matching responder URL. \\
R08--R09 & Resume interrupted work without extra copies and retain actionable
status information for partially completed entries. \\
R10 & Protect private artifacts and document cloud access and temporary-resource
cleanup. Use fictional names in committed test fixtures. \\
\hline
\end{tabularx}
\end{center}

\subsection{Approval and Privacy}
The control spreadsheet is the working surface; the generated roster is a
separate deliverable. Cancellation before approval must not create an output
roster or student forms. Temporary OCR resources may already exist and must
follow a documented cleanup policy on success, failure, and cancellation.

Real screenshots, rosters, credentials, and exported run state must not enter
version control. Local ignore rules support this policy but do not replace
review of staged files. Cloud permissions, retention, and cleanup mechanisms
remain design decisions to validate. A form bearing a student's name does not
authenticate the respondent; access settings must be tested separately.

